\documentclass[10pt]{article}
\usepackage[T1]{fontenc}
\usepackage[utf8]{inputenc}
\usepackage[french]{babel}
\usepackage{amsmath,amssymb,mathrsfs,mathpazo}
\usepackage{geometry}
\geometry{a4paper,margin=2cm}
\usepackage{graphicx,array,tabularx,tabularray,colortbl,multirow,multicol,booktabs,enumitem,pifont,ulem,fancybox,diagbox,cancel}
\usepackage{pstricks,pstricks-add,pst-plot,pst-tree,pst-node,pst-text,pst-func,pst-math,pst-bspline,pst-3dplot,pst-eucl,pst-solides3d}
\let\pscalculate\undefined
\usepackage{pst-bezier}
\usepackage{tikz,pgfplots,tkz-tab}
\usetikzlibrary{arrows,arrows.meta,patterns,calc,positioning,trees,decorations.markings,shapes.geometric,shapes.misc,shapes.arrows,intersections,decorations.pathmorphing,decorations.pathreplacing,backgrounds,shadows,babel}
\pgfplotsset{compat=1.18}
\usepgfplotslibrary{fillbetween}
\usepackage[np]{numprint}
\usepackage{listings,makecell,pas-tableur,twoopt}
\let\barre\undefined
\usepackage[pstricks]{tablvar}
\definecolor{darkgreen}{RGB}{0,100,0}
\providecommand{\red}{\color{red}}
\providecommand{\blue}{\color{blue}}
\providecommand{\green}{\color{green}}
\providecommand{\black}{\color{black}}
\providecommand{\white}{\color{white}}
\providecommand{\gray}{\color{gray}}
\providecommand{\cyan}{\color{cyan}}
\providecommand{\magenta}{\color{magenta}}
\providecommand{\yellow}{\color{yellow}}
\providecommand{\violet}{\color{violet}}
\providecommand{\orange}{\color{orange}}
\providecommand{\euro}{\texteuro}
\providecommand{\EUR}{\texteuro}
\providecommand{\ieme}{\textsuperscript{e}}
\providecommand{\degres}{\ensuremath{^\circ}}
\providecommand{\degree}{\ensuremath{^\circ}}
\providecommand{\vv}[1]{\overrightarrow{#1}}
\providecommand{\Coord}[1]{#1}
\providecommand{\decosix}{}
\providecommand{\starredbullet}{\ensuremath{\bullet}}
\pagestyle{empty}
\setlength{\parindent}{0pt}
\usepackage{adjustbox}\DeclareUnicodeCharacter{2236}{:}\DeclareUnicodeCharacter{2190}{\ensuremath{\leftarrow}}
\providecommand{\R}{\mathbb{R}}
\providecommand{\N}{\mathbb{N}}
\providecommand{\Z}{\mathbb{Z}}
\providecommand{\Q}{\mathbb{Q}}
\providecommand{\D}{\mathbb{D}}
\providecommand{\C}{\mathbb{C}}
\providecommand{\e}{{\mathrm{e}}}
\providecommand{\vect}[1]{\overrightarrow{#1}}
\providecommand{\vectt}[1]{\overrightarrow{\mathrm{#1}}}
\providecommand{\vvt}[1]{\overrightarrow{\mathrm{#1}}}
\providecommand{\np}[2][]{{#2~#1}}
\providecommand{\vv}[1]{{\overrightarrow{#1}}}
\renewcommand{\d}{{\mathrm{d}}}
\providecommand{\degres}{{^\circ}}
\providecommand{\degre}{{^\circ}}
\providecommand{\blue}{\color{blue}}
\providecommand{\red}{\color{red}}
\providecommand{\enskip}{\quad}
\providecommand{\llbracket}{\lbrack\!\lbrack}
\providecommand{\rrbracket}{\rbrack\!\rbrack}
\providecommand{\euro}{{\text{€}}}
\providecommand{\no}{{\text{n°}}}
\providecommand{\interval}[2]{{[#1;#2]}}
\providecommand{\intervFO}[2]{{[#1;#2[}}
\providecommand{\dotfill}{\ldots}

\providecommand{\ds}{\displaystyle}
\providecommand{\textding}[1]{#1}
\providecommand{\Oij}{\(\mathrm{O};\vec{\imath},\vec{\jmath}\)}
\providecommand{\Oijk}{\(\mathrm{O};\vec{\imath},\vec{\jmath},\vec{k}\)}
\providecommand{\Ouv}{\(\mathrm{O};\vec{u},\vec{v}\)}
\providecommand{\intervFF}[2]{\ensuremath{[#1;#2]}}
\renewcommand{\euro}{\ifmmode\text{\texteuro}\else\texteuro\fi}
\DeclareUnicodeCharacter{00B0}{\ifmmode^{\circ}\else\textdegree\fi}
\DeclareUnicodeCharacter{20AC}{\ifmmode\text{\texteuro}\else\texteuro\fi}
% Énoncés : Éducation nationale. Transcription et corrigés : APMEP, Denis Vergès et contributeurs. Découpage et classement : CoopMaths, Sébastien Lozano, Éric Elter et contributeurs. Adaptation de la présentation : Maths Explorer.
% CC BY-SA : https://coopmaths.fr/www/about/
% https://coopmaths.fr/alea/static/bac/2024/tex/bac_2024_06_sujet2_asie_4.tex
\begin{document}
\noindent\textit{APMEP / CoopMaths. CC BY-SA. Présentation adaptée par Maths Explorer.}\par\medskip
\medskip

Dans un repère orthonormé \Oijk{} de l'espace, on considère le plan $(P)$ d'équation: 
\[(P) :\quad 2x + 2y - 3z + 1 = 0.\]

On considère les trois points A, B et C de coordonnées: 
\begin{center}A(1~;~0~;~1) , B$(2~;~- 1~;~1)$ et C$( - 4~;~- 6~;~5)$ .\end{center}

Le but de cet exercice est d'étudier le rapport des aires entre un triangle et son projeté orthogonal dans un plan.


\medskip

\textbf{Partie A}

\begin{enumerate}
\item Pour chacun des points A, B et C, vérifier s'il appartient au plan $(P)$.
\item Montrer que le point C$'(0~;~- 2~;~- 1)$ est le projeté orthogonal du point C sur le plan $(P)$.
\item Déterminer une représentation paramétrique de la droite (AB).
\item On admet l'existence d'un unique point H vérifiant les deux conditions 

$\left\{\begin{array}{l}
\text{H} \in \text{(AB)}\\
\text{(AB) et (HC) sont orthogonales}.
\end{array}\right.$

Déterminer les coordonnées du point H.

\begin{center}
\psset{unit=1cm}
\begin{pspicture}(7.8,3.3)
\pspolygon(0,0)(5.5,0)(7.8,3.2)(2.3,3.2)
\pspolygon(2.3,0.6)(4.7,1.8)(2.6,2.8)
\uput[d](2.3,0.6){A} \uput[r](4.7,1.8){B} \uput[ur](2.6,2.8){C}
\psline[linestyle=dashed](2.3,0.6)(2.6,1.8)(2.6,2.8)
\psline[linestyle=dashed](2.6,1.8)(4.7,1.8)
\uput[ur](2.6,1.8){C$'$}\rput(5.4,0.2){$(P)$}\end{pspicture}
\end{center}

\end{enumerate}

\medskip

\textbf{Partie B}

\medskip

On admet que les coordonnées du vecteur $\vect{\text{HC}}$ sont : $\vect{\text{HC}}\begin{pmatrix}- \frac{11}{2}\\- \frac{11}{2}\\4\end{pmatrix}$.

\medskip

\begin{enumerate}
\item Calculer la valeur exacte de $\left\|\vect{\text{HC}}\right\|$.
\item Soit $S$ l'aire du triangle ABC. Déterminer la valeur exacte de $S$.
\end{enumerate}

\medskip

\textbf{Partie C}

\medskip

On admet que HC$'= \sqrt{\dfrac{17}{2}}$.

\medskip

\begin{enumerate}
\item Soit $\alpha = \widehat{\text{CHC}'}$. Déterminer la valeur de $\cos(\alpha)$.
\item 
	\begin{enumerate}
		\item Montrer que les droites (C$'$H) et (AB) sont perpendiculaires.
		\item Calculer $S'$ l'aire du triangle ABC$'$, on donnera la valeur exacte.
		\item Donner une relation entre $S,\: S'$ et $\cos (\alpha)$.
	\end{enumerate}
\end{enumerate}

\end{document}
